% META: {"id": "1NSI-WEB-IHM-RE-C9", "chapitre": "1NSI-WEB-IHM", "type_objet": "remediation", "capacites": ["P-IHM-04C"], "status": "generated"}

%---------------------------------------------------------------
% FICHE DE REMEDIATION — C9 : ne pas utiliser GET pour des donnees confidentielles
%---------------------------------------------------------------

\begin{exercice}{1NSI-WEB-RE-C9-EX1}{1}{10}
Un élève conçoit un formulaire de réinitialisation de mot de passe avec la méthode GET :
\begin{console}
<form action="/nouveau_mdp" method="GET">
    <input type="password" name="nouveau_mdp">
    <button type="submit">Valider</button>
</form>
\end{console}
\begin{enumerate}
  \item En utilisant \lstinline{urlencode}, construire l'URL qui serait générée si l'utilisateur saisit \lstinline{"MonSecret2026"} comme nouveau mot de passe.
  \item Le nouveau mot de passe apparaît-il dans cette URL ? Où cette URL risque-t-elle d'être enregistrée (historique, journaux serveur) ?
  \item Réécrire l'attribut \lstinline{method} du formulaire pour corriger ce problème.
\end{enumerate}
\end{exercice}

% BEGIN-VERIFY
% from urllib.parse import urlencode
% donnees = {"nouveau_mdp": "MonSecret2026"}
% url = "https://site.fr/nouveau_mdp?" + urlencode(donnees)
% assert "MonSecret2026" in url
% END-VERIFY

\begin{corrige}{1NSI-WEB-RE-C9-EX1}
\textbf{1.}
\begin{python}
from urllib.parse import urlencode

donnees = {"nouveau_mdp": "MonSecret2026"}
url = "https://site.fr/nouveau_mdp?" + urlencode(donnees)
print(url)
# https://site.fr/nouveau_mdp?nouveau_mdp=MonSecret2026
\end{python}

\textbf{2.} Oui, le mot de passe apparaît \textbf{en clair} dans l'URL. Cette URL risque
d'être enregistrée dans l'\textbf{historique du navigateur} et dans les
\textbf{journaux (logs) du serveur}, deux endroits où un mot de passe ne devrait jamais
apparaître, même sur une connexion HTTPS (qui protège seulement la transmission, pas ces
enregistrements).

\textbf{3.}
\begin{console}
<form action="/nouveau_mdp" method="POST">
    <input type="password" name="nouveau_mdp">
    <button type="submit">Valider</button>
</form>
\end{console}
Avec \lstinline{method="POST"}, le mot de passe est placé dans le corps de la requête, pas
dans l'URL.
\end{corrige}
